% Auto-generated from raw.txt
% Usage:
%   \vocabentry{word}{/ipa/ (pos)}{definition}{example}{note}

\vocabentry{Astronaut}
{/ˈæstrənɔːt/ (n)}
{A person who is trained to travel in a spacecraft.}
{Becoming an astronaut requires years of intensive physical and mental training.}
{Đây là một danh từ dùng để chỉ phi hành gia, người được đào tạo để làm việc trong không gian.}

\vocabentry{Cosmos}
{/ˈkɑːzmoʊs/ (n)}
{The universe seen as a well-ordered whole.}
{Ancient philosophers spent their lives trying to understand the mysteries of the cosmos.}
{Đây là một danh từ dùng để chỉ vũ trụ, thường mang sắc thái bao la và có trật tự.}

\vocabentry{Spacecraft}
{/ˈspeɪskræft/ (n)}
{A vehicle or machine designed to fly in outer space.}
{The unmanned spacecraft successfully landed on the surface of Mars last night.}
{Đây là một danh từ dùng để chỉ tàu vũ trụ hoặc các thiết bị bay trong không gian.}

\vocabentry{Orbit}
{/ˈɔːrbɪt/ (v)}
{To move in a curved path around a planet, moon, or star.}
{Thousands of pieces of space junk are currently orbiting the Earth at high speeds.}
{Đây là một động từ dùng để chỉ việc quay quanh quỹ đạo.}

\vocabentry{Celestial}
{/səˈlestʃl/ (adj)}
{Positioned in or relating to the sky, or outer space as observed in astronomy.}
{Stars, planets, and comets are all examples of celestial bodies.}
{Đây là một tính từ dùng để chỉ các thiên thể hoặc những gì thuộc về bầu trời/vũ trụ.}

\vocabentry{Gravity}
{/ˈɡrævəti/ (n)}
{The force that attracts a body toward the center of the earth, or toward any other physical body having mass.}
{Astronauts experience weightlessness because of the lack of gravity in outer space.}
{Đây là một danh từ dùng để chỉ trọng lực.}

\vocabentry{Satellite}
{/ˈsætəlaɪt/ (n)}
{An artificial body placed in orbit around the earth or moon or another planet in order to collect information or for communication.}
{Weather satellites provide essential data for predicting hurricanes and storms.}
{Đây là một danh từ dùng để chỉ vệ tinh nhân tạo hoặc vệ tinh tự nhiên.}

\vocabentry{Interstellar}
{/ˌɪntərˈstelər/ (adj)}
{Occurring or situated between stars.}
{Interstellar travel remains a dream for scientists due to the vast distances involved.}
{Đây là một tính từ dùng để chỉ không gian hoặc hành trình giữa các vì sao.}

\vocabentry{Galaxy}
{/ˈɡæləksi/ (n)}
{A system of millions or billions of stars, together with gas and dust, held together by gravitational attraction.}
{The Milky Way is the galaxy that contains our entire solar system.}
{Đây là một danh từ dùng để chỉ thiên hà.}

\vocabentry{Telescope}
{/ˈtelɪskoʊp/ (n)}
{An optical instrument designed to make distant objects appear nearer.}
{The James Webb Space Telescope has captured stunning images of early galaxies.}
{Đây là một danh từ dùng để chỉ kính thiên văn.}

\vocabentry{Rocketry}
{/ˈrɑːkɪtri/ (n)}
{The branch of science and technology concerned with rockets.}
{Modern rocketry has allowed humans to send probes to the outer edges of the solar system.}
{Đây là một danh từ dùng để chỉ ngành chế tạo và sử dụng tên lửa.}

\vocabentry{Launch}
{/lɔːntʃ/ (v/n)}
{To send a rocket, missile, or satellite into the air or upper atmosphere. (Đây là một động từ/danh từ chỉ việc phóng (tên lửa/tàu vũ trụ) hoặc buổi phóng.).}
{Thousands of people gathered to watch the launch of the latest lunar mission.}
{}

\vocabentry{Colonization}
{/ˌkɑːlənəˈzeɪʃn/ (n)}
{The action or process of settling among and establishing control over a place, often a planet in this context. (Đây là một danh từ dùng để chỉ sự định cư hoặc thuộc địa hóa (ví dụ: định cư trên sao Hỏa).).}
{Elon Musk has often spoken about the eventual colonization of Mars to ensure human survival.}
{}

\vocabentry{Atmosphere}
{/ˈætməsfɪr/ (n)}
{The envelope of gases surrounding the earth or another planet.}
{Mars has a very thin atmosphere compared to Earth, making it difficult for humans to breathe.}
{Đây là một danh từ dùng để chỉ bầu khí quyển.}

\vocabentry{Extraterrestrial}
{/ˌekstrətəˈrestriəl/ (adj)}
{Of or from outside the earth or its atmosphere.}
{Scientists are constantly searching for signs of extraterrestrial life on distant moons.}
{Đây là một tính từ dùng để chỉ những gì ngoài Trái Đất hoặc thuộc về người ngoài hành tinh.}

\vocabentry{Probe}
{/proʊb/ (n)}
{An unmanned exploratory spacecraft that transmits information from its surroundings.}
{The Voyager probe has been traveling through space for over forty years.}
{Đây là một danh từ dùng để chỉ tàu thăm dò không người lái.}

\vocabentry{Solar System}
{/ˈsoʊlər ˌsɪstəm/ (n)}
{The collection of eight planets and their moons in orbit around the sun.}
{There are eight confirmed planets in our solar system, with Earth being the third from the Sun.}
{Đây là một cụm danh từ dùng để chỉ Hệ Mặt Trời.}

\vocabentry{Black hole}
{/ˈblæk hoʊl/ (n)}
{A region of space having a gravitational field so intense that no matter or radiation can escape.}
{A black hole is formed when a massive star collapses at the end of its life cycle.}
{Đây là một cụm danh từ dùng để chỉ hố đen.}

\vocabentry{Astrophysics}
{/ˌæstroʊˈfɪzɪks/ (n)}
{The branch of astronomy concerned with the physical nature of stars and other celestial bodies.}
{Studying astrophysics helps us understand how the universe began and how it will end.}
{Đây là một danh từ dùng để chỉ ngành vật lý thiên văn.}

\vocabentry{Nebula}
{/ˈnebjələ/ (n)}
{A cloud of gas and dust in outer space, visible in the night sky as an indistinct bright patch.}
{Stars are born within vast clouds of gas known as a nebula.}
{Đây là một danh từ dùng để chỉ tinh vân.}

\vocabentry{Eclipse}
{/ɪˈklɪps/ (v)}
{(Of a celestial body) obscure the light from or to (another celestial body).}
{The Moon will eclipse the Sun for several minutes during the afternoon.}
{Đây là một động từ dùng để chỉ việc che khuất ánh sáng.}

\vocabentry{Meteor}
{/ˈmiːtiər/ (n)}
{A small body of matter from outer space that enters the earth's atmosphere, becoming incandescent as a result of friction.}
{We saw a bright meteor streak across the sky during the meteor shower last night.}
{Đây là một danh từ dùng để chỉ sao băng.}

\vocabentry{Comet}
{/ˈkɑːmət/ (n)}
{A celestial object consisting of a nucleus of ice and dust and, when near the sun, a ‘tail’ of gas and dust particles pointing away from the sun.}
{Halley's Comet is visible from Earth only once every 75 to 76 years.}
{Đây là một danh từ dùng để chỉ sao chổi.}

\vocabentry{Supernova}
{/ˌsuːpərˈnoʊvə/ (n)}
{A star that suddenly increases greatly in brightness because of a catastrophic explosion that ejects most of its mass.}
{A supernova explosion can briefly outshine an entire galaxy.}
{Đây là một danh từ dùng để chỉ siêu tân tinh.}

\vocabentry{Vacuum}
{/ˈvækjuːm/ (n)}
{A space entirely devoid of matter.}
{Sound cannot travel through the vacuum of space because there are no air molecules to carry it.}
{Đây là một danh từ dùng để chỉ môi trường chân không trong không gian.}

\vocabentry{Radiation}
{/ˌreɪdiˈeɪʃn/ (n)}
{The emission of energy as electromagnetic waves or as moving subatomic particles.}
{Spacecraft must be heavily shielded to protect crews from harmful cosmic radiation.}
{Đây là một danh từ dùng để chỉ bức xạ vũ trụ, một mối nguy hiểm cho phi hành gia.}

\vocabentry{Trajectory}
{/trəˈdʒektəri/ (n)}
{The path followed by a moving object through the air or in space.}
{A small error in the spacecraft's trajectory could cause it to miss its target planet.}
{Đây là một danh từ dùng để chỉ quỹ đạo/đường bay.}

\vocabentry{Thruster}
{/ˈθrʌstər/ (n)}
{A small rocket engine used for maneuvering or maintaining the attitude of a spacecraft.}
{The pilot used the thrusters to dock the capsule with the space station.}
{Đây là một danh từ dùng để chỉ động cơ đẩy điều hướng.}

\vocabentry{Payload}
{/ˈpeɪloʊd/ (n)}
{The cargo carried by a vehicle, especially by a spacecraft or rocket.}
{The cost of a space mission is largely determined by the weight of the payload.}
{Đây là một danh từ dùng để chỉ tải trọng hữu ích.}

\vocabentry{Re-entry}
{/riː ˈentri/ (n)}
{The action of an object re-entering the earth's atmosphere from outer space.}
{During re-entry, the heat shield must withstand temperatures of thousands of degrees.}
{Đây là một danh từ dùng để chỉ quá trình tàu vũ trụ trở lại bầu khí quyển Trái Đất.}

\vocabentry{Exoplanet}
{/ˈeksoʊˌplænɪt/ (n)}
{A planet which orbits a star outside the solar system.}
{Astronomers have found thousands of exoplanets that might have liquid water.}
{Đây là một danh từ dùng để chỉ hành tinh ngoài hệ mặt trời.}

\vocabentry{Propulsion}
{/prəˈpʌlʃn/ (n)}
{The action of driving or pushing forward.}
{Ion propulsion is a highly efficient technology used for long-distance space missions.}
{Đây là một danh từ dùng để chỉ hệ thống đẩy hoặc lực đẩy của tên lửa.}

\vocabentry{Capsule}
{/ˈkæpsuːl/ (n)}
{A small spacecraft or part of a larger one that is designed to return to Earth with a crew or instruments.}
{The astronauts returned safely to Earth in a small metal capsule.}
{Đây là một danh từ dùng để chỉ khoang tàu vũ trụ.}

\vocabentry{Module}
{/ˈmɑːdʒuːl/ (n)}
{Each of a set of parts or units that can be used to construct a more complex structure, such as a space station.}
{The Russian service module provides power and life support to the space station.}
{Đây là một danh từ dùng để chỉ các khoang hoặc mô-đun lắp ghép.}

\vocabentry{Lunar}
{/ˈluːnər/ (adj)}
{Of, determined by, or resembling the moon.}
{The Apollo missions provided us with a wealth of information about the lunar surface.}
{Đây là một tính từ dùng để chỉ những gì thuộc về mặt trăng.}

\vocabentry{Altitude}
{/ˈæltɪtuːd/ (n)}
{The height of an object or point in relation to sea level or ground level.}
{The satellite is orbiting at an altitude of five hundred kilometers.}
{Đây là một danh từ dùng để chỉ độ cao so với mặt nước biển hoặc bề mặt hành tinh.}

\vocabentry{Asteroid}
{/ˈæstərɔɪd/ (n)}
{A small rocky body orbiting the sun.}
{NASA is developing a plan to deflect any asteroid that might threaten the Earth.}
{Đây là một danh từ dùng để chỉ tiểu hành tinh.}

\vocabentry{Observatory}
{/əbˈzɜːrvətɔːri/ (n)}
{A room or building housing an astronomical telescope or other scientific equipment for the study of natural phenomena.}
{The observatory is located on top of a mountain to avoid light pollution from the city.}
{Đây là một danh từ dùng để chỉ đài quan sát thiên văn.}

\vocabentry{Cosmology}
{/kɑːzˈmɑːlədʒi/ (n)}
{The science of the origin and development of the universe.}
{Modern cosmology is based on the Big Bang theory.}
{Đây là một danh từ dùng để chỉ vũ trụ học.}

\vocabentry{Docking}
{/ˈdɑːkɪŋ/ (n)}
{The process of two spacecraft joining together in space.}
{The docking procedure was successful, allowing the crew to move between the two ships.}
{Đây là một danh từ dùng để chỉ việc kết nối/ghép nối các tàu vũ trụ trong không gian.}

\vocabentry{Cosmonaut}
{/ˈkɑːzmənɔːt/ (n)}
{A Russian astronaut.}
{Yuri Gagarin was the first cosmonaut to orbit the Earth in 1961.}
{Đây là một danh từ dùng để chỉ phi hành gia người Nga.}

\vocabentry{Telemetry}
{/təˈlemətri/ (n)}
{The automatic measurement and wireless transmission of data from remote sources.}
{The scientists used telemetry to monitor the health of the astronauts during their mission.}
{Đây là một danh từ dùng để chỉ hệ thống đo xa/truyền dữ liệu từ xa.}

\vocabentry{Light-year}
{/ˈlaɪt jɪr/ (n)}
{A unit of astronomical distance equivalent to the distance that light travels in one year.}
{The nearest star system, Alpha Centauri, is about 4.3 light-years away.}
{Đây là một danh từ dùng để chỉ năm ánh sáng.}

\vocabentry{Spectroscopy}
{/spekˈtrɑːskəpi/ (n)}
{The branch of science concerned with the investigation and measurement of spectra produced when matter interacts with or emits electromagnetic radiation.}
{Spectroscopy allows astronomers to determine the chemical composition of distant stars.}
{Đây là một danh từ dùng để chỉ quang phổ học, giúp xác định thành phần của các ngôi sao.}

\vocabentry{Parallax}
{/ˈpærəlæks/ (n)}
{The effect whereby the position or direction of an object appears to differ when viewed from different positions.}
{Stellar parallax is used to measure the distance of nearby stars from our solar system.}
{Đây là một danh từ dùng để chỉ thị sai, phương pháp đo khoảng cách các vì sao.}

\vocabentry{Manifest}
{/ˈmænɪfest/ (n)}
{A document listing the cargo, passengers, and crew of a ship or aircraft.}
{The mission manifest included several scientific experiments and two new satellites.}
{Trong không gian, đây là danh sách hàng hóa/thiết bị mang theo.}

\vocabentry{Avionics}
{/ˌeɪviˈɑːnɪks/ (n)}
{Electronic equipment used in aircraft and spacecraft.}
{A failure in the ship's avionics made it difficult for the pilot to maintain control.}
{Đây là một danh từ dùng để chỉ hệ thống điện tử hàng không vũ trụ.}

\vocabentry{Centrifuge}
{/ˈsentrɪfjuːdʒ/ (n)}
{A machine with a rapidly rotating container that applies centrifugal force to its contents.}
{Astronauts train in a large centrifuge to prepare for the intense forces of launch and re-entry.}
{Dùng để huấn luyện phi hành gia chịu lực G.}

\vocabentry{Inertia}
{/ɪˈnɜːrʃə/ (n)}
{The tendency of an object to resist changes in its motion (to stay at rest or keep moving).}
{Due to inertia, an astronaut floating in space will continue to drift unless they use a thruster.}
{Đây là một danh từ dùng để chỉ quán tính.}

\vocabentry{Combustion}
{/kəmˈbʌstʃən/ (n)}
{The process of burning something.}
{Efficient combustion in the rocket engine is necessary to produce maximum thrust.}
{Trong tên lửa, đây là sự đốt cháy nhiên liệu.}

\vocabentry{Cryogenic}
{/ˌkraɪoʊˈdʒenɪk/ (adj)}
{Relating to or involving the production of very low temperatures.}
{Many modern rockets use cryogenic liquid oxygen and hydrogen as fuel.}
{Đây là một tính từ dùng để chỉ nhiên liệu tên lửa siêu lạnh.}

\vocabentry{Hypervelocity}
{/ˌhaɪpərvəˈlɑːsəti/ (n)}
{Extremely high velocity.}
{Space debris traveling at hypervelocity can cause significant damage to satellites.}
{Đây là một danh từ dùng để chỉ vận tốc siêu cực lớn.}

\vocabentry{Orbiter}
{/ˈɔːrbɪtər/ (n)}
{A spacecraft designed to go into orbit around a celestial body without landing on it.}
{The Mars Reconnaissance Orbiter has been studying the planet's surface for many years.}
{Đây là một danh từ dùng để chỉ tàu quay quanh quỹ đạo.}

\vocabentry{Lander}
{/ˈlændər/ (n)}
{A spacecraft designed to land on the surface of a planet or moon.}
{The lunar lander carried the first humans to the surface of the Moon in 1969.}
{Đây là một danh từ dùng để chỉ tàu đổ bộ.}

\vocabentry{Rover}
{/ˈroʊvər/ (n)}
{A space exploration vehicle designed to move across the surface of a planet or other celestial body.}
{The Curiosity rover has discovered evidence of ancient water on Mars.}
{Đây là một danh từ dùng để chỉ xe tự hành thăm dò bề mặt hành tinh.}

\vocabentry{Aeronautics}
{/ˌerəˈnɔːtɪks/ (n)}
{The science or practice of travel through the air.}
{NASA stands for the National Aeronautics and Space Administration.}
{Đây là một danh từ dùng để chỉ ngành hàng không.}

\vocabentry{Simulated}
{/ˈsɪmjuleɪtɪd/ (adj)}
{Manufactured in imitation of some other object; mock.}
{Astronauts spend hundreds of hours in simulated space environments before their mission.}
{Dùng để chỉ các môi trường giả lập huấn luyện.}

\vocabentry{Microgravity}
{/ˌmaɪkroʊˈɡrævəti/ (n)}
{Very weak gravity, as in an orbiting spacecraft.}
{Experiments conducted in microgravity provide unique insights into physical processes.}
{Đây là một danh từ dùng để chỉ môi trường vi trọng lực.}

\vocabentry{Life-support}
{/ˈlaɪf səˌpɔːrt/ (n)}
{Equipment that allows a person to survive in a place where they would normally die.}
{A failure in the life-support system is one of the most dangerous situations in space.}
{Đây là một cụm danh từ dùng để chỉ hệ thống duy trì sự sống.}

\vocabentry{Space suit}
{/ˈspeɪs suːt/ (n)}
{A garment designed to allow an astronaut to survive in space.}
{A space suit provides oxygen, temperature control, and protection from radiation.}
{Đây là một cụm danh từ dùng để chỉ bộ đồ vũ trụ.}

\vocabentry{Spacewalk}
{/ˈspeɪswɔːk/ (n)}
{An act of physical activity done by an astronaut in space outside a spacecraft.}
{The astronauts performed a six-hour spacewalk to repair a broken antenna.}
{Đây là một danh từ dùng để chỉ việc đi bộ ngoài không gian.}

\vocabentry{Astrobiology}
{/ˌæstroʊbaɪˈɑːlə dʒi/ (n)}
{The study of the origin, evolution, and possibility of life in the universe.}
{Astrobiology is a multidisciplinary field that combines biology, chemistry, and physics.}
{Đây là một danh từ dùng để chỉ ngành sinh học vũ trụ.}

\vocabentry{Radio telescope}
{/ˈreɪdioʊ ˈtelɪskoʊp/ (n)}
{An instrument used to detect radio emissions from sky objects.}
{Radio telescopes can see objects that are invisible to optical telescopes.}
{Đây là một cụm danh từ dùng để chỉ kính thiên văn vô tuyến.}

\vocabentry{Interplanetary}
{/ˌɪntərˈplænəteri/ (adj)}
{Situated or occurring between planets.}
{Interplanetary missions require complex calculations to account for planetary motion.}
{Đây là một tính từ dùng để chỉ các hành trình giữa các hành tinh.}

\vocabentry{Horizon}
{/həˈraɪzn/ (n)}
{The line at which the earth's surface and the sky appear to meet; in space, the "event horizon" of a black hole.}
{Nothing, not even light, can escape from beyond the event horizon of a black hole.}
{Đây là một danh từ dùng để chỉ đường chân trời hoặc chân trời sự kiện.}

\vocabentry{Magnetosphere}
{/mæɡˈniːtəsfɪr/ (n)}
{The region surrounding the earth or another astronomical body in which its magnetic field is the predominant effective magnetic field.}
{The Earth's magnetosphere protects us from the solar wind and cosmic rays.}
{Đây là một danh từ dùng để chỉ từ quyển.}

\vocabentry{Solar wind}
{/ˈsoʊlər wɪnd/ (n)}
{A continuous stream of charged particles from the sun.}
{The solar wind can interfere with satellite communications and power grids on Earth.}
{Đây là một cụm danh từ dùng để chỉ gió mặt trời.}

\vocabentry{Quasar}
{/ˈkweɪzɑːr/ (n)}
{A massive and extremely remote celestial object, emitting exceptionally large amounts of energy.}
{Quasars are among the brightest and most distant objects in the universe.}
{Đây là một danh từ dùng để chỉ chuẩn tinh.}

\vocabentry{Pulsar}
{/ˈpʌlsɑːr/ (n)}
{A celestial object, thought to be a rapidly rotating neutron star, that emits regular pulses of radio waves.}
{The regular pulses of a pulsar can be used as extremely accurate cosmic clocks.}
{Đây là một danh từ dùng để chỉ sao xung.}

\vocabentry{Dark matter}
{/dɑːrk ˈmætər/ (n)}
{Non-luminous material that is postulated to exist in space.}
{Scientists believe that dark matter makes up about 85\% of the matter in the universe.}
{Đây là một cụm danh từ dùng để chỉ vật chất tối.}

\vocabentry{Dark energy}
{/dɑːrk ˈenərdʒi/ (n)}
{A theoretical repulsive force that counteracts gravity and causes the universe to expand at an accelerating rate.}
{Dark energy is thought to be responsible for the accelerating expansion of the universe.}
{Đây là một cụm danh từ dùng để chỉ năng lượng tối.}

\vocabentry{Singularity}
{/ˌsɪŋɡjuˈlærəti/ (n)}
{A point at which a function takes an infinite value, especially in space-time at the center of a black hole.}
{All the mass of a black hole is concentrated in a single point called a singularity.}
{Đây là một danh từ dùng để chỉ điểm kỳ dị.}

\vocabentry{Spacetime}
{/ˈspeɪstaɪm/ (n)}
{The four-dimensional continuum of the three spatial coordinates plus time.}
{Einstein's theory of general relativity describes gravity as the curvature of spacetime.}
{Đây là một danh từ dùng để chỉ không-thời gian.}

\vocabentry{Equilibrium}
{/ˌiːkwɪˈlɪbriəm/ (n)}
{A state in which opposing forces or influences are balanced.}
{A star remains stable as long as there is an equilibrium between gravity and nuclear fusion.}
{Trong vũ trụ học là trạng thái cân bằng của các ngôi sao.}

\vocabentry{Accretion}
{/əˈkriːʃn/ (n)}
{The coming together and cohesion of matter under the influence of gravitation to form larger bodies.}
{Planets are formed through the accretion of dust and gas in a young solar system.}
{Đây là một danh từ dùng để chỉ sự bồi tụ vật chất.}

\vocabentry{Escape velocity}
{/ɪˈskeɪp vəˌlɑːsəti/ (n)}
{The lowest velocity which a body must have in order to escape the gravitational attraction of a particular planet or other object.}
{A rocket must reach escape velocity to break free from Earth's gravity and travel to other planets.}
{Đây là một cụm danh từ dùng để chỉ vận tốc thoát.}

\vocabentry{Solar flare}
{/ˈsoʊlər fler/ (n)}
{A brief eruption of intense high-energy radiation from the sun's surface.}
{A massive solar flare can disrupt electronics on Earth and pose a risk to astronauts.}
{Đây là một cụm danh từ dùng để chỉ bùng nổ mặt trời.}

\vocabentry{Sunspot}
{/ˈsʌnspɑːt/ (n)}
{A spot or patch appearing from time to time on the sun's surface, appearing dark by contrast with its surroundings.}
{Sunspots are regions of reduced surface temperature caused by intense magnetic activity.}
{Đây là một danh từ dùng để chỉ vết đen mặt trời.}

\vocabentry{Astrochemistry}
{/ˌæstroʊˈkemɪstri/ (n)}
{The branch of science concerned with the study of the chemical substances and species occurring in stars and interstellar space.}
{Astrochemistry helps us understand how organic molecules can form in space.}
{Đây là một danh từ dùng để chỉ ngành hóa học vũ trụ.}

\vocabentry{Geocentric}
{/ˌdʒiːoʊˈsentrɪk/ (adj)}
{Having or representing the earth as the center. (Đây là một tính từ dùng để chỉ thuyết địa tâm (Trái Đất là trung tâm).).}
{The geocentric model of the universe was replaced by the heliocentric model centuries ago.}
{}

\vocabentry{Heliocentric}
{/ˌhiːlioʊˈsentrɪk/ (adj)}
{Having or representing the sun as the center. (Đây là một tính từ dùng để chỉ thuyết nhật tâm (Mặt Trời là trung tâm).).}
{Copernicus was one of the first scientists to propose a heliocentric model of our solar system.}
{}

\vocabentry{Apogee}
{/ˈæpədʒiː/ (n)}
{The point in the orbit of the moon or a satellite at which it is furthest from the earth. (Đây là một danh từ dùng để chỉ điểm viễn địa (điểm xa Trái Đất nhất trên quỹ đạo).).}
{The satellite reaches its apogee at a distance of thirty-six thousand kilometers.}
{}

\vocabentry{Perigee}
{/ˈperɪdʒiː/ (n)}
{The point in the orbit of the moon or a satellite at which it is nearest to the earth. (Đây là một danh từ dùng để chỉ điểm cận địa (điểm gần Trái Đất nhất trên quỹ đạo).).}
{At perigee, the Moon appears slightly larger and brighter in the night sky.}
{}

\vocabentry{Perihelion}
{/ˌperɪˈhiːliən/ (n)}
{The point in the orbit of a planet, comet, or other body at which it is closest to the sun.}
{The Earth reaches perihelion in early January each year.}
{Đây là một danh từ dùng để chỉ điểm cận nhật.}

\vocabentry{Aphelion}
{/æpˈhiːliən/ (n)}
{The point in the orbit of a planet, comet, or other body at which it is furthest from the sun.}
{Despite being at aphelion in July, the Northern Hemisphere is warmer due to the Earth's tilt.}
{Đây là một danh từ dùng để chỉ điểm viễn nhật.}

\vocabentry{Nadir}
{/ˈneɪdɪr/ (n)}
{The point on the celestial sphere directly below an observer.}
{In astronomy, the nadir is the point directly opposite the zenith.}
{Đây là một danh từ dùng để chỉ điểm thiên để, đối diện với thiên đỉnh.}

\vocabentry{Zenith}
{/ˈziːnɪθ/ (n)}
{The point on the celestial sphere vertically above an observer.}
{The sun reaches its zenith at midday in the tropics.}
{Đây là một danh từ dùng để chỉ thiên đỉnh.}

\vocabentry{Declination}
{/ˌdeklɪˈneɪʃn/ (n)}
{The angular distance of a point north or south of the celestial equator. (Đây là một danh từ dùng để chỉ xích vĩ (tọa độ trên bầu trời).).}
{Declination and right ascension are used to map the positions of stars.}
{}

\vocabentry{Ascension}
{/əˈsenʃn/ (n)}
{The act of rising to an important position or a higher level.}
{Astronomers use right ascension to pinpoint the exact location of a celestial object.}
{Trong thiên văn là "Right Ascension" - xích kinh.}

\vocabentry{Magnitude}
{/ˈmæɡnɪtuːd/ (n)}
{The great size or importance of something; in astronomy, the brightness of a star.}
{Stars with a lower magnitude value are actually brighter than those with higher values.}
{Đây là một danh từ dùng để chỉ cấp sao hoặc độ sáng của một thiên thể.}

\vocabentry{Parsec}
{/ˈpɑːrsek/ (n)}
{A unit of distance used in astronomy, equal to about 3.26 light-years.}
{Most distant galaxies are located millions or even billions of parsecs away.}
{Đây là một danh từ dùng để chỉ đơn vị thị sai giây - parsec.}

\vocabentry{Light-pollution}
{/laɪt pəˈluːʃn/ (n)}
{Brightening of the night sky caused by street lights and other man-made sources.}
{Light pollution makes it very difficult for city dwellers to see the stars at night.}
{Đây là một cụm danh từ dùng để chỉ ô nhiễm ánh sáng.}

\vocabentry{Constellation}
{/ˌkɑːnstəˈleɪʃn/ (n)}
{A group of stars forming a recognizable pattern.}
{Orion is one of the most easily recognized constellations in the winter sky.}
{Đây là một danh từ dùng để chỉ chòm sao.}

\vocabentry{Zodiac}
{/ˈzoʊdiæk/ (n)}
{A belt of the heavens within about 8° either side of the ecliptic.}
{The twelve signs of the zodiac are based on the constellations the sun passes through.}
{Đây là một danh từ dùng để chỉ cung hoàng đạo.}

\vocabentry{Equinox}
{/ˈekwɪnɑːks/ (n)}
{The time or date (twice each year) at which the sun crosses the celestial equator. (Đây là một danh từ dùng để chỉ điểm phân (xuân phân hoặc thu phân).).}
{During the equinox, day and night are of approximately equal length all over the world.}
{}

\vocabentry{Solstice}
{/ˈsɑːlstɪs/ (n)}
{Either of the two times in the year when the sun reaches its highest or lowest point in the sky at noon. (Đây là một danh từ dùng để chỉ điểm chí (hạ chí hoặc đông chí).).}
{The summer solstice is the longest day of the year in the Northern Hemisphere.}
{}

\vocabentry{Rotate}
{/ˈroʊteɪt/ (v)}
{Move or cause to move in a circle around an axis or center.}
{The Earth rotates on its axis once every twenty-four hours.}
{Đây là một động từ dùng để chỉ sự tự quay quanh trục.}

\vocabentry{Revolution}
{/ˌrevəˈluːʃn/ (n)}
{The movement of one object around a center or another object.}
{The Earth's revolution around the Sun causes the change of seasons.}
{Đây là một danh từ dùng để chỉ sự chuyển động trên quỹ đạo.}

\vocabentry{Tidal}
{/ˈtaɪdl/ (adj)}
{Relating to or affected by tides.}
{Tidal forces from the Moon cause the oceans on Earth to rise and fall.}
{Tính từ chỉ thuộc về thủy triều hoặc lực triều.}

\vocabentry{Crater}
{/ˈkreɪtər/ (n)}
{A large, bowl-shaped cavity in the ground or on the surface of a planet or the moon.}
{The surface of the Moon is covered in thousands of impact craters.}
{Đây là một danh từ dùng để chỉ hố va chạm hoặc miệng núi lửa.}

\vocabentry{Mantle}
{/ˈmæntl/ (n)}
{The region of the interior of a terrestrial planet between the core and the crust.}
{The Earth's mantle is made of solid rock that flows very slowly over time.}
{Đây là một danh từ dùng để chỉ lớp manti của hành tinh.}

\vocabentry{Core}
{/kɔːr/ (n)}
{The central or innermost portion of the earth or another planet.}
{The Earth's outer core is made of liquid iron and nickel.}
{Đây là một danh từ dùng để chỉ lõi hành tinh hoặc nhân sao.}

\vocabentry{Crust}
{/krʌst/ (n)}
{The outermost layer of a terrestrial planet.}
{The Earth's crust is divided into several large tectonic plates.}
{Đây là một danh từ dùng để chỉ lớp vỏ hành tinh.}

\vocabentry{Basalt}
{/bəˈsɔːlt/ (n)}
{A dark, fine-grained volcanic rock that sometimes makes up planetary surfaces.}
{Large areas of the Moon are covered in basaltic plains known as "maria."}
{Đây là một danh từ dùng để chỉ đá bazan.}

\vocabentry{Regolith}
{/ˈreɡəlɪθ/ (n)}
{A layer of loose, heterogeneous superficial deposits covering solid rock.}
{Lunar regolith is a fine dust created by billions of years of meteorite impacts.}
{Đây là một danh từ chuyên ngành chỉ lớp bụi đất trên bề mặt mặt trăng/hành tinh.}

\vocabentry{Ecliptic}
{/ɪˈklɪptɪk/ (n)}
{A great circle on the celestial sphere representing the sun's apparent path during the year.}
{Most of the planets in our solar system orbit near the plane of the ecliptic.}
{Đây là một danh từ dùng để chỉ hoàng đạo.}

\vocabentry{Inclination}
{/ˌɪnklɪˈneɪʃn/ (n)}
{The angle between the orbital plane of a planet and the ecliptic.}
{Pluto has a much higher orbital inclination than the eight major planets.}
{Đây là một danh từ dùng để chỉ độ nghiêng quỹ đạo.}

\vocabentry{Retrograde}
{/ˈretrəɡreɪd/ (adj)}
{(Of a planet) moving in the opposite direction to that of other bodies in its system.}
{Venus is unusual because it has a retrograde rotation, spinning clockwise on its axis.}
{Đây là một tính từ dùng để chỉ sự chuyển động nghịch hành.}

\vocabentry{Geosynchronous}
{/ˌdʒiːoʊˈsɪŋkrənəs/ (adj)}
{(Of an earth satellite) having a period of rotation of 24 hours.}
{Communications satellites are often placed in a geosynchronous orbit to remain above the same spot on Earth.}
{Đây là một tính từ dùng để chỉ quỹ đạo địa đồng bộ.}

\vocabentry{Geostationary}
{/ˌdʒiːoʊˈsteɪʃəneri/ (adj)}
{(Of an earth satellite) orbiting so that it remains over the same point on the earth's surface.}
{A geostationary satellite appears to stay in the same position in the sky at all times.}
{Đây là một tính từ dùng để chỉ quỹ đạo địa tĩnh.}

\vocabentry{Aerospace}
{/ˈeroʊspeɪs/ (n)}
{The branch of technology and industry concerned with both aviation and space flight.}
{He decided to pursue a career in aerospace engineering to design next-generation rockets.}
{Đây là một danh từ dùng để chỉ ngành hàng không vũ trụ.}

\vocabentry{Astrography}
{/əˈstrɑːɡrəfi/ (n)}
{The branch of astronomy that deals with the mapping of stars and celestial bodies.}
{Astrography has become much more accurate thanks to digital imaging and space telescopes.}
{Đây là một danh từ dùng để chỉ thiên văn đồ bản học.}

\vocabentry{Astrometry}
{/əˈstrɑːmətri/ (n)}
{The branch of astronomy that involves precise measurements of the positions and movements of stars.}
{Astrometry is essential for detecting the subtle wobble of stars caused by orbiting planets.}
{Đây là một danh từ dùng để chỉ thiên văn trắc địa.}

\vocabentry{Astronomy}
{/əˈstrɑːnəmi/ (n)}
{The branch of science that deals with celestial objects, space, and the physical universe as a whole.}
{Astronomy is one of the oldest sciences, dating back to ancient civilizations.}
{Đây là một danh từ dùng để chỉ ngành thiên văn học.}

\vocabentry{Bolometer}
{/boʊˈlɑːmɪtər/ (n)}
{An instrument for measuring the power of incident electromagnetic radiation.}
{Bolometers are used to detect faint infrared radiation from distant stars and galaxies.}
{Đây là một danh từ chỉ nhiệt kế bức xạ dùng trong thiên văn.}

\vocabentry{Binary star}
{/ˈbaɪneri stɑːr/ (n)}
{A system of two stars in which one revolves around the other or both revolve around a common center.}
{About half of all the stars in our galaxy are thought to be part of a binary star system.}
{Đây là một cụm danh từ dùng để chỉ sao đôi.}

\vocabentry{Neutron star}
{/ˈnuːtrɑːn stɑːr/ (n)}
{A celestial object of very small radius and very high density composed predominantly of closely packed neutrons.}
{A neutron star is so dense that a single teaspoon of its material would weigh billions of tons.}
{Đây là một cụm danh từ dùng để chỉ sao neutron.}

\vocabentry{White dwarf}
{/waɪt dwɔːrf/ (n)}
{A small very dense star that is typically the size of a planet.}
{Our Sun will eventually become a white dwarf at the end of its life cycle.}
{Đây là một cụm danh từ dùng để chỉ sao lùn trắng.}

\vocabentry{Red giant}
{/red ˈdʒaɪənt/ (n)}
{A very large star of high luminosity and relatively low surface temperature.}
{When the Sun runs out of hydrogen fuel, it will expand into a red giant and swallow the inner planets.}
{Đây là một cụm danh từ dùng để chỉ sao khổng lồ đỏ.}

\vocabentry{Main sequence}
{/meɪn ˈsiːkwəns/ (n)}
{A series of star types to which most stars belong.}
{The Sun is currently in the main sequence stage of its evolution.}
{Đây là một cụm danh từ dùng để chỉ dãy chính trong sơ đồ tiến hóa sao.}

\vocabentry{Interstellar medium}
{/ˌɪntərˈstelər ˈmiːdiəm/ (n)}
{The matter that exists in the space between the star systems in a galaxy.}
{The interstellar medium consists of gas, dust, and cosmic rays.}
{Đây là một cụm danh từ dùng để chỉ môi trường liên sao.}

\vocabentry{Molecular cloud}
{/məˈlekjələr klaʊd/ (n)}
{An interstellar cloud of gas and dust in which molecules can form.}
{Stars are formed when regions within a molecular cloud collapse under their own gravity.}
{Đây là một cụm danh từ dùng để chỉ đám mây phân tử, nơi sinh ra các ngôi sao.}

\vocabentry{Planetary nebula}
{/ˌplænəteri ˈnebjələ/ (n)}
{A ring-shaped nebula formed by an expanding shell of gas around an aging star.}
{Despite their name, planetary nebulae have nothing to do with planets.}
{Đây là một cụm danh từ dùng để chỉ tinh vân hành tinh.}

\vocabentry{Impact}
{/ˈɪmpækt/ (n/v)}
{The action of one object coming forcibly into contact with another.}
{A massive asteroid impact is thought to have caused the extinction of the dinosaurs.}
{Trong không gian, đây là sự va chạm thiên thể.}

\vocabentry{Meteoroid}
{/ˈmiːtiərɔɪd/ (n)}
{A small body moving in the solar system that would become a meteor if it entered the earth's atmosphere.}
{Most meteoroids are fragments of asteroids or comets.}
{Đây là một danh từ dùng để chỉ thiên thạch khi còn ở trong không gian.}

\vocabentry{Meteorite}
{/ˈmiːtiəraɪt/ (n)}
{A meteor that survives its passage through the earth's atmosphere such that part of it strikes the ground.}
{Scientists study meteorites to learn about the early history of the solar system.}
{Đây là một danh từ dùng để chỉ mảnh thiên thạch đã rơi xuống mặt đất.}

\vocabentry{Meteor shower}
{/ˈmiːtiər ˈʃaʊər/ (n)}
{A number of meteors that appear to radiate from one point in the sky.}
{The Perseid meteor shower is one of the most popular astronomical events of the year.}
{Đây là một cụm danh từ dùng để chỉ mưa sao băng.}

\vocabentry{Propellant}
{/prəˈpelənt/ (n)}
{A substance used to provide thrust in a rocket engine.}
{A rocket must carry both fuel and an oxidizer as part of its propellant system.}
{Đây là một danh từ dùng để chỉ nhiên liệu tên lửa/chất đẩy.}

\vocabentry{Radioactivity}
{/ˌreɪdioʊækˈtɪvəti/ (n)}
{The emission of ionizing radiation or particles.}
{Cosmic rays are a form of high-energy radioactivity that comes from outside our solar system.}
{Trong không gian, đây là hiện tượng phóng xạ tự nhiên.}

\vocabentry{Shielding}
{/ˈʃiːldɪŋ/ (n)}
{Material used to protect someone or something from a danger, especially radiation.}
{Lead and water are two materials that can be used for effective radiation shielding.}
{Đây là một danh từ dùng để chỉ lớp chắn bảo vệ.}

\vocabentry{Symmetry}
{/ˈsɪmətri/ (n)}
{The quality of being made up of exactly similar parts facing each other.}
{Aerodynamic symmetry is essential for rockets that must travel through the Earth's atmosphere.}
{Trong thiết kế tàu vũ trụ là sự cân đối.}

\vocabentry{Thermodynamics}
{/ˌθɜːrmoʊdaɪˈnæmɪks/ (n)}
{The branch of physical science that deals with the relations between heat and other forms of energy.}
{Understanding thermodynamics is essential for designing efficient rocket engines.}
{Trong tên lửa học là nhiệt động lực học.}

\vocabentry{Thrust}
{/θrʌst/ (n/v)}
{The propulsive force of a jet or rocket engine.}
{The rocket engine produces millions of pounds of thrust to lift the vehicle off the launch pad.}
{Đây là một danh từ/động từ dùng để chỉ lực đẩy.}

\vocabentry{Titanium}
{/taɪˈteɪniəm/ (n)}
{A strong, low-density, highly corrosion-resistant lustrous transition metal.}
{Titanium is widely used in the aerospace industry because it is both strong and lightweight.}
{Vật liệu chính chế tạo tàu vũ trụ.}

\vocabentry{Total}
{/ˈtoʊtl/ (adj)}
{Complete; including everything.}
{Thousands of people traveled to the path of totality to witness the total solar eclipse.}
{Ví dụ: "total eclipse" - nhật thực toàn phần.}

\vocabentry{Transmission}
{/trænzˈmɪʃn/ (n)}
{The action or process of transmitting something.}
{It takes about twenty minutes for a radio transmission to travel from Mars to Earth.}
{Sự truyền tín hiệu từ không gian về Trái Đất.}

\vocabentry{Transparency}
{/trænsˈpærənsi/ (n)}
{The condition of being transparent.}
{Good atmospheric transparency is essential for high-quality astronomical observations.}
{Trong thiên văn là độ trong của bầu khí quyển.}

\vocabentry{Trend}
{/trend/ (n)}
{A general direction in which something is developing or changing.}
{The current trend in space exploration is the increasing involvement of private companies like SpaceX.}
{Trong khám phá không gian là xu hướng tư nhân hóa.}

\vocabentry{Trial}
{/ˈtraɪəl/ (n)}
{A formal examination of evidence; or a test of performance.}
{The new rocket engine underwent several static fire trials before its first flight.}
{Các đợt thử nghiệm thiết bị vũ trụ.}

\vocabentry{Trough}
{/trɔːf/ (n)}
{A long, narrow open container; or a low point in a wave.}
{Astronomers study the peaks and troughs of light waves to understand the properties of stars.}
{Trong sóng vô tuyến vũ trụ.}

\vocabentry{Turbulent}
{/ˈtɜːrbjələnt/ (adj)}
{Characterized by conflict, disorder, or confusion; not controlled or calm.}
{The spacecraft experienced a turbulent ride as it passed through the upper atmosphere.}
{Tính chất của khí quyển hoặc gió mặt trời.}

\vocabentry{Typical}
{/ˈtɪpɪkl/ (adj)}
{Having the distinctive qualities of a particular type of person or thing.}
{A typical spiral galaxy contains hundreds of billions of stars and large amounts of gas and dust.}
{Ví dụ: "typical galaxy" - thiên hà điển hình.}

\vocabentry{Ubiquitous}
{/juːˈbɪkwɪtəs/ (adj)}
{Present, appearing, or found everywhere.}
{Hydrogen is the most ubiquitous element in the entire universe.}
{Trong vũ trụ là hydrogen.}

\vocabentry{Ultimate}
{/ˈʌltɪmət/ (adj)}
{Being or happening at the end of a process; final.}
{The ultimate goal of many space agencies is to send humans to another star system.}
{Mục tiêu cuối cùng của khám phá không gian.}

\vocabentry{Underground}
{/ˌʌndərˈɡraʊnd/ (adj)}
{Beneath the surface of the ground.}
{Scientists are looking for evidence of underground liquid water on Jupiter's moon, Europa.}
{Ví dụ: "underground water" trên các hành tinh khác.}

\vocabentry{Underlying}
{/ˌʌndərˈlaɪɪŋ/ (adj)}
{Significant as a cause or basis of something but not necessarily manifest or obvious.}
{General relativity is the underlying theory that explains how gravity works on a cosmic scale.}
{Lý thuyết nền tảng.}

\vocabentry{Undertake}
{/ˌʌndərˈteɪk/ (v)}
{Commit oneself to and begin (an enterprise or responsibility).}
{Several countries are planning to undertake a joint mission to explore the outer planets.}
{Bắt đầu một nhiệm vụ vũ trụ.}

\vocabentry{Unfavorable}
{/ʌnˈfeɪvərəbl/ (adj)}
{Expressing or showing a lack of approval or support.}
{The launch was postponed due to unfavorable weather conditions at the space center.}
{Ví dụ: "unfavorable weather" cho buổi phóng tàu.}

\vocabentry{Uniqueness}
{/juˈniːknəs/ (n)}
{The quality of being particularly remarkable, special, or unusual.}
{Space exploration has highlighted the uniqueness and fragility of our home planet.}
{Tính độc nhất của Trái Đất.}

\vocabentry{Universal}
{/ˌjuːnɪˈvɜːrsl/ (adj)}
{Of, affecting, or done by all people or things in the world or in a particular group.}
{The laws of physics are believed to be universal throughout the entire cosmos.}
{Tính phổ quát của các định luật vật lý.}

\vocabentry{Unlimited}
{/ʌnˈlɪmɪtɪd/ (adj)}
{Not limited or restricted in terms of number, quantity, or extent.}
{Outer space offers virtually unlimited resources for future human generations.}
{Ví dụ: "unlimited potential" của không gian.}

\vocabentry{Unprecedented}
{/ʌnˈpresɪdentɪd/ (adj)}
{Never done or known before.}
{The successful landing on a moving comet was an unprecedented scientific achievement.}
{Những thành tựu chưa từng có tiền lệ.}

\vocabentry{Unregulated}
{/ʌnˈreɡjuleɪtɪd/ (adj)}
{Not controlled or supervised by regulations or laws.}
{There are concerns about the unregulated growth of the commercial space industry.}
{Ví dụ: "unregulated space commercialization".}

\vocabentry{Unstable}
{/ʌnˈsteɪbl/ (adj)}
{Prone to change, fail, or give way; not stable.}
{Some massive stars become unstable at the end of their lives and explode as supernovae.}
{Các ngôi sao không ổn định.}

\vocabentry{Untapped}
{/ˌʌnˈtæpt/ (adj)}
{(Of a resource) not yet exploited or used.}
{The Moon may contain vast untapped reserves of helium-3 and water ice.}
{Các nguồn tài nguyên chưa được khai phá trên mặt trăng/tiểu hành tinh.}

\vocabentry{Utility}
{/juːˈtɪləti/ (n)}
{The state of being useful, profitable, or beneficial.}
{The utility of satellites for global communications and navigation is immense.}
{Lợi ích thực tế của vệ tinh.}

\vocabentry{Utilization}
{/ˌjuːtələˈzeɪʃn/ (n)}
{The action of making practical and effective use of something.}
{The utilization of Martian soil to create building materials is a key goal for future missions.}
{Sự sử dụng tài nguyên tại chỗ - ISRU.}

\vocabentry{Valuable}
{/ˈvæljuəbl/ (adj)}
{Worth a great deal of money.}
{Asteroid mining could provide a new source of valuable metals like platinum and gold.}
{Các khoáng sản quý trên tiểu hành tinh.}

\vocabentry{Value}
{/ˈvæljuː/ (n)}
{The importance, worth, or usefulness of something.}
{The scientific value of the samples returned from the Moon is immeasurable.}
{Giá trị khoa học của nhiệm vụ.}

\vocabentry{Vanish}
{/ˈvænɪʃ/ (v)}
{Disappear suddenly and completely.}
{Any object that crosses the event horizon of a black hole will vanish forever.}
{Thiên thể biến mất vào hố đen.}

\vocabentry{Variable}
{/ˈveriəbl/ (adj)}
{Not consistent or having a fixed pattern; liable to change.}
{A variable star is one whose brightness changes over time, often in a regular cycle.}
{Ví dụ: "variable star" - sao biến quang.}

\vocabentry{Vast}
{/væst/ (adj)}
{Of very great extent or quantity; immense.}
{The distance between galaxies is so vast that it is difficult for the human mind to grasp.}
{Sự bao la của không gian.}

\vocabentry{Vegetation}
{/ˌvedʒəˈteɪʃn/ (n)}
{Plants considered collectively.}
{Scientists are looking for chemical signatures that might indicate the presence of vegetation on exoplanets.}
{Trong sinh học vũ trụ, tìm kiếm dấu hiệu thực vật.}

\vocabentry{Vehicle}
{/ˈviːəkl/ (n)}
{A thing used for transporting people or goods.}
{The Space Shuttle was a reusable launch vehicle that served NASA for thirty years.}
{Trong không gian là phương tiện vận chuyển - Launch Vehicle.}

\vocabentry{Velocity}
{/vəˈlɑːsəti/ (n)}
{The speed of something in a given direction.}
{A spacecraft must maintain a specific orbital velocity to avoid falling back to Earth.}
{Vận tốc.}

\vocabentry{Venture}
{/ˈventʃər/ (n/v)}
{A risky or daring journey or undertaking.}
{Human spaceflight remains a dangerous and expensive venture for any nation.}
{Chuyến thám hiểm mạo hiểm vào không gian.}

\vocabentry{Versatile}
{/ˈvɜːrsətl/ (adj)}
{Able to adapt or be adapted to many different functions or activities.}
{The robotic arm on the space station is a versatile tool used for many different tasks.}
{Thiết bị vũ trụ đa năng.}

\vocabentry{Vessel}
{/ˈvesl/ (n)}
{A ship or large boat.}
{The alien vessel appeared on the radar screens of the military base.}
{Trong không gian là tàu vũ trụ.}

\vocabentry{Viability}
{/ˌvaɪəˈbɪləti/ (n)}
{Ability to work successfully.}
{Scientists are studying the viability of growing crops in Martian soil.}
{Tính khả thi của cuộc sống trên sao Hỏa.}

\vocabentry{Vibrant}
{/ˈvaɪbrənt/ (adj)}
{Full of energy and life.}
{The Hubble Space Telescope has captured many vibrant images of colorful nebulae.}
{Hình ảnh rực rỡ của các thiên hà.}

\vocabentry{Vibration}
{/vaɪˈbreɪʃn/ (n)}
{An instance of vibrating.}
{Spacecraft components must be tested to ensure they can withstand the intense vibrations of launch.}
{Rung lắc khi phóng tên lửa.}

\vocabentry{Vigor}
{/ˈvɪɡər/ (n)}
{Physical strength and good health.}
{Astronauts must maintain their physical vigor through regular exercise in weightlessness.}
{Sức khỏe của phi hành gia.}

\vocabentry{Virtual}
{/ˈvɜːrtʃuəl/ (adj)}
{Almost or nearly as described, but not completely.}
{Virtual reality is used to train astronauts for complex repairs on the space station.}
{Thực tế ảo trong huấn luyện.}

\vocabentry{Visibility}
{/ˌvɪzəˈbɪləti/ (n)}
{The state of being able to see or be seen.}
{The visibility of the planet Mars increases significantly during its closest approach to Earth.}
{Độ hiển thị của các vì sao.}

\vocabentry{Vital}
{/ˈvaɪtl/ (adj)}
{Absolutely necessary or important; essential.}
{Oxygen and water are vital for the survival of any human crew in space.}
{Cần thiết cho sự sống.}

\vocabentry{Volatility}
{/ˌvɑːləˈtɪləti/ (n)}
{Liability to change rapidly and unpredictably.}
{The volatility of water ice on the Moon's surface depends on its temperature and location.}
{Tính dễ bay hơi của các chất trên hành tinh.}

\vocabentry{Volume}
{/ˈvɑːljuːm/ (n)}
{The amount of space that a substance or object occupies.}
{Jupiter is so large that its volume is more than one thousand times that of the Earth.}
{Thể tích của một ngôi sao/hành tinh.}

\vocabentry{Voluntary}
{/ˈvɑːlənteri/ (adj)}
{Done, given, or acting of one's own free will.}
{Citizen science projects rely on the voluntary contributions of amateur astronomers around the world.}
{Sự tham gia tự nguyện của cộng đồng vào thiên văn học.}

\vocabentry{Vulnerability}
{/ˌvʌlnərəˈbɪləti/ (n)}
{The quality or state of being exposed to the possibility of being attacked or harmed.}
{The vulnerability of spacecraft to micrometeoroids is a constant concern for mission planners.}
{Sự dễ bị tổn thương của tàu vũ trụ trước vi thiên thạch.}

\vocabentry{Ward}
{/wɔːrd/ (v)}
{Guard and protect.}
{Special magnetic fields could be used to ward off harmful cosmic rays from a space colony.}
{Trong không gian là việc bảo vệ khỏi bức xạ.}

\vocabentry{Waste}
{/weɪst/ (n)}
{Material which is eliminated or discarded as no longer useful.}
{Space debris, also known as space waste, is becoming a serious problem for satellites in Earth's orbit.}
{Rác thải không gian - Space Junk.}

\vocabentry{Wavelength}
{/ˈweɪvleŋθ/ (n)}
{The distance between successive crests of a wave.}
{Different telescopes are designed to observe the universe at different wavelengths of light.}
{Bước sóng ánh sáng/vô tuyến.}

\vocabentry{Wealth}
{/welθ/ (n)}
{An abundance of valuable possessions or money.}
{Space missions have provided us with a wealth of information about the history of our solar system.}
{Ví dụ: "wealth of information" - lượng thông tin phong phú.}

\vocabentry{Weather}
{/ˈweðər/ (n)}
{The state of the atmosphere at a place and time.}
{Space weather is driven by activity on the Sun and can affect satellites and power grids.}
{Trong không gian là "space weather" - thời tiết không gian.}

\vocabentry{Weigh}
{/weɪ/ (v)}
{Find out how heavy (someone or something) is.}
{It is difficult to weigh objects in space, so scientists measure their mass instead.}
{Đo khối lượng trong môi trường không trọng lực.}

\vocabentry{Weightless}
{/ˈwaɪtləs/ (adj)}
{(Of a person or object) appearing to have no weight.}
{Astronauts have fun floating and doing flips in the weightless environment of the space station.}
{Tình trạng không trọng lượng.}

\vocabentry{Western}
{/ˈwestərn/ (adj)}
{Living in or originating from the West.}
{Western space agencies like NASA and ESA collaborate on many international projects.}
{Các cơ quan vũ trụ phương Tây.}

\vocabentry{Widespread}
{/ˌwaɪdˈspred/ (adj)}
{Found or distributed over a large area or number of people.}
{There is widespread public interest in the possibility of discovering life on other planets.}
{Sự quan tâm rộng rãi của công chúng.}

\vocabentry{Wild}
{/waɪld/ (adj)}
{(Of an animal or plant) living or growing in the natural environment.}
{There are many wild theories about what lies at the center of a black hole.}
{Ví dụ: "wild theories" - những giả thuyết táo bạo.}

\vocabentry{Zone}
{/zoʊn/ (n)}
{An area or stretch of land having a particular characteristic.}
{The habitable zone is the region around a star where liquid water can exist on a planet's surface.}
{Ví dụ: "Habitable Zone" - vùng ở được quanh một ngôi sao.}

